\chapter{LITERATURE REVIEW}

\section{Introduction}

The following section delves into published research that is relevant to the design and construction of TerraNode. The review begins by examining the overall current context of digital agriculture and then moves to an exploration of the different elements of blockchain technology, such as blockchain basics, food supply chain traceability, predictive crop analytics and the characteristics of the Sui blockchain platform. Limitations of the current approaches are highlighted and an explicit research statement is made for each section. This project to encompass that gap was identified. Finally, a summary table has been provided for comparison so that the shortcomings of the available solution, and TerraNode's differences, can be discussed.

\section{Overview of Precision Agriculture and Digital Farming}

Precision Agriculture is a new invention in response to growing food demands and land scarcity in facing weather variability and increased food production. It requires sensor network, satellite photography, as well as computing models to track the crops and soil and weather at a very micro level. The fundamental is to take decisions about the extent to which a farmer should manage their farm according to real-time data rather than tradition or rough estimation.

\textcite{padhiary2025emerging} provided a comprehensive survey on the technologies which enable smart and sustainable precision agriculture. They explored the field of remote sensing, robotic systems, UAVs, and Internet of Things (IoT) sensor platforms. They concluded that although single technologies have come a long way, they have yet to progress into a more comprehensive and integrated solution. Bundling components into a single decision-support platform, still remains a major challenge. One of the major findings was that most of the existing implementations consider data collection and data analysis as two independent processes. By itself, the effect on the agility of farmers' ability to act on the insights at the right time is restricted.

The modes of development, the technologies used and analysed security issues discussed on smart agriculture is also analysed by \textcite{raj2022survey}. Thus, while numerous deployments are currently underway with a set of monitoring domotics based on IoT devices, their use in centralized cloud architectures poses risks of single points of failures and enables collection of sensitive data for loss or compromise, their analysis reveals. They said that the next generation of information-based agricultural systems would need to handle data integrity and computational security processes on par with the other technologies, and add them in as a protective measure.

The transition from traditional farm records to the use of sensors in farming has not varied from region to region. Though agriculture has come a long way in the last few decades, there are still smallholder farmers with less than 20 acres who rely on paper logs and experience to run their business who live in many parts of the developing world. This gap between technology adopted and technology available is the background which motivates the design of TerraNode with respect to how a technically effective and user friendly system will be built for different types of users.

\section{Agricultural Data Management Systems}

One of the foundations of any digital farming endeavor is agricultural data management systems. They collect environmental readings, store the data, then report this data to the project, present summarized data to decision-makers, and keep historical production records. For the most part, these systems have been established in a relational database model with a centralized server.

In the systematic literature review, \textcite{sharma2023blockchain} studied the literature related to a food supply chain management system. They reviewed the existing management systems available and identified that the majority of the systems were fragile, with data silos, manual paper work and data interoperability problems across stakeholders identified. From over 300 published studies they conducted a bibliometric analysis and revealed that even though there is considerable interest in blockchain for supply chain management since 2018, the majority of systems are developed at a conceptual or proof-of-concept stage that don't progress to a level that can be field-tested to meet the needs of smallholder farmers in developing countries.

One riddle with centralized agricultural databases is that they are vulnerable to data manipulation. Having to trust the single ownership of the database could result in alteration of the historical data for any reason, including fraud or administrative error. This drawback may be of great annoyance when the data is to be transferred to other systems to be used for developing analytical models. Where the records do not demonstrate authenticity of originals, then any projections, or recommendations for such projections, are necessarily incorrect. This is solved by TerraNode by linking each telemetry record and produce batch to Sui blockchain before any kind of processing is performed, and so you could mathematically confirm the data origin before processing.

\section{Blockchain Technology Fundamentals}

Blockchain technology offers an append-only data structure that is decentralized. Records are arranged in blocks which are consecutively cryptographically linked and propagated from one-node to another via a node network. After the network consensus confirms a block, the information carried becomes almost unchangeable. This attribute is what makes blockchain a popular choice for applications which need auditability, transparency, and trust.

The researchers of \textcite{chauhan2025blockchain} had mentioned the link between blockchain and the big data analytics in the agriculture and also about the possibility of implementing digital ledger technology in agriculture and using massive amounts of data to create intelligent agricultural systems. They categorized the different types of blockchain consensus mechanisms into proof-of-work, proof-of-stake and delegated proof-of-stake types and revealed that these schemes require much less energy consumption and attain much greater transaction speed that is vital in an agriculture context, where hardware is limited.

Blockchains may be implemented through various architectures and each of them has its advantages and disadvantages; the most decentralized blockchain is the public blockchain type such as Bitcoin or Ethereum, but computational resources are intensive. Developers should expect a capped number of transactions per second and fees will be more expensive. But private blockchains like Hyperledger Fabric get better performance for the same downsides: loss of trustless properties that makes public blockchains appealing. In the middle is a Consortium blockchain that has a handful of validators that manage the network. The Sui blockchain TerraNode is built on is public, has a delegated proof-of-stake (dPoS) consensus mechanism, delivers high throughput, low transaction fees, and continues to be transparent and open like a public network.


\section{Blockchain Applications in Agricultural Supply Chains}

The application of blockchain technology in the agricultural supply chain can be done in various ways. The use of blockchain technology in agricultural supply chains has been a subject of great interest over the last few years, and has been a focus of research. The central theme is that blockchain can be used to create an immutable history of all transactions and custody movements that take place from the farm gate to the end consumer, thus allowing for full traceability and minimising opportunities for fraud.

\textcite{caro2018blockchain} introduced AgriBlockIoT, a fully decentralized blockchain-based traceability solution for agri-food supply chain management. Their system was designed to integrate devices generating and consuming digital data on the supply chain when implemented as IoT. They used two different blockchain solutions, one on Ethereum and one with Hyperledger Sawtooth, which was measured by latency, number of CPU, memory, and network overhead. The deployment of Hyperledger was quicker because it was proof-of-work based, but the deployments of Ethereum were slower since they were more decentralised. This balance of performance and trustlessness appears throughout the literature and is a key issue addressed by the delegated proof-of-stake consensus mechanism in the Sui blockchain.

\textcite{mirabelli2020blockchain} studied the relationship between blockchain and agricultural supply chain traceability by means of a structured literature review. Their analysis identified and discussed key research trends such as the increasing adoption of smart contracts for automating compliance checks, the incorporation of IoT sensor data into on-chain records, and the difficulty of scaling blockchain systems to accommodate the large amount of data produced by large agricultural operations. They said that though there are a number of proof-of-concept implementations, there are not many that have been tested under realistic conditions in the field with real farmers as end users.

\textcite{rahman2025enhancing} explored the transparency of supply chain in a data-driven approach in the context of distributed ledger applications. Their research showed that trust between supply chain stakeholders was significantly higher in a blockchain-based supply chain than in conventional supply chains. But they also discovered some barriers to adoption, including the energy needed to power mining consensus and a lack of uniformity in how agricultural information is recorded on-chain.

\textcite{lin2017blockchain} proposed an early and significant argument for blockchain as the next step in Information and Communications Technology (ICT) for agriculture. They presented their work from an environmental systems engineering point of view and suggested that by introducing blockchain to the agricultural sector, baseline environmental information would remain intact for all stakeholders involved. Their work informed many of the design decisions for TerraNode, for example that they want to take environmental data and record it on an unalterable chain of blocks before they use it for downstream analysis.

\section{Blockchain-Based Smart Farm Security Frameworks}

Apart from tracing the flow of goods, blockchain technology has been proposed as a foundation for complete farm security systems that protect not just the information of the supply chain but the operational technology systems running in today's agriculture.

Aliyu et al.~\parencite{aliyu2023blockchain} proposed a smart farm security framework using the Internet of Things (IoT) on the blockchain platform as an alternative solution. The issue of weak control in IoT sensor/actuator communications with the centralized management server is resolved in their system. They have a system to store sensor data and control instructions onto a blockchain, ensuring that alterations of a device's configuration are not possible to conceal. They have a close connection to data integrity as with TerraNode, but their efforts are geared towards system protection as they lack predictive analytics capabilities as well as a supply chain custody tracking module.

\textcite{agtrust2026framework} presented the AgTrust system to safeguard critical agriculture technology and new solutions. It is a large framework of guidelines that establishes a comprehensive agriculture technology security solution from the edge and field to the network and analytics in the cloud. AgTrust provided a threat taxonomy for agricultural IoT environments and recommended mitigation measures for each layer. The framework gave useful security ideas, however, it did not create a viable prototype or show how to incorporate with a particular blockchain platform.

\textcite{dzreke2025securing} took a closer look at the connected farm cybersecurity topic by experimenting on the vulnerabilities of IoT devices. The study uncovered 32 commercial agricultural IoT devices and 15 farm subsystems with attack surfaces. The results demonstrated that it was possible to make up to 19 percent changes to crop prescriptions, pointing to the direct link between cyber intrusions and biological results. This paper emphasizes the importance of data integrity tools, including the ones implemented in TerraNode - SHA-256 hashing, blockchain anchoring, and more - which ensure that data cannot be altered by unauthorized parties.

\section{Predictive Analytics and Crop Yield Forecasting}

Predictive analytics is a new area of interest in agricultural informatics. Good prediction of yields helps farmers to make decisions about resource allocation, harvest timing and marketing. This problem has been tackled by many statistical and machine learning methods.

\textcite{khaki2019cnn} proposed a deep learning model that combines Convolutional Network (CNN) and Recurrent Network (RNN) for crop yield prediction. Their model was trained using environmental data and management practices across the US Corn Belt and found to have a root-mean-square-error of 9 percent for corn yield and 8 percent for soybean yield, outperforming models using the random forest, deep fully-connected neural network and LASSO regression. They showed how temporal dependencies in environmental data can be modeled using deep learning models to inform the development of more sophisticated versions of the TerraNode analytics engine.

\textcite{garg2018crop} studied crop yield prediction using fuzzy logic and regression models. They used fuzzy time series analysis to the wheat yield data for the past, making experiments on different partitioning schemes of intervals and different regression equations. They found that fuzzy logic-based models were able to cope with the ambiguity of agricultural information better than the deterministic models. Their model is based on the regression concept, similar to what TerraNode uses with its current analytics engine, which uses telemetry to extract environmental features rather than the yield timeseries from the past.

\textcite{layek2025secured} have introduced a secured trifecta of IoT, machine learning and blockchain in the field of crop forecasting. Their solution combined data collection from the IoT with machine learning prediction models with data verification through blockchain. It is one of the conceptually closest pieces to TerraNode, with the understanding that data security is key to the reliability of predictions. However, their system was not implemented on a public ledger, thus restricting the benefits of transparency and verifiability. Besides, they lacked what is required for a real-world multi-stakeholder application: a web-based management interface and role-based access control (RBAC).

\section{Supply Chain Traceability and Distributed Ledger Systems}

The process of tracing the origin of agricultural products from the farm to the consumer is a complex task which involves a number of actors with different rights and responsibilities. A distributed ledger technology can provide a natural solution to the problem of recording a custody transfer transparently among all parties and that's less likely to be modified retroactively.

\textcite{cocco2021blockchain} presented the use of blockchain and self-sovereign identity to enable the food supply chain to be supported in terms of quality assurance. Their system enabled individual actors to maintain their own identity credentials while also being part of a common product quality attestation ledger. The self-sovereign identity concept is also very inspirational for the wallet login system of TerraNode, which allows users to log in with their Sui wallet address and a cryptographic signature, rather than just centralized credentials.

\textcite{sun2025storage} tackled a real-world issue related to blockchain technology in agriculture: the storage overhead for full-node replication. Their Storage Light Node (SLN) model featured a classification of cold/hot data based on identity relevance, generation time and query frequency. They demonstrated that 50,563 real agricultural supply chain records showed a 95.10 per cent reduction in storage usage for full-node operation. The paper focuses on scalability issues, which will be increasingly important as TerraNode transitions from prototype to product.

\textcite{ktari2022agricultural} introduced a lightweight embedded blockchain system for agricultural supply chain traceability, using the case study of olive oil production in Tunisia. They were connected to the IoT sensors and ran directly on embedded hardware at the production site. Their light-weight block chain was able to be run on resource-constrained farms common in developing nations. Although TerraNode doesn't focus on embedded deployments, the value of minimizing computational overhead is consistent with the decision to use the Sui blockchain, which has reduced costs and is less expensive than Ethereum-based solutions when it comes to transaction costs.

\section{The Sui Blockchain and Move Programming Language}

Built by Mysten Labs, the Sui blockchain is a next-generation layer-one blockchain that employs a delegated proof-of-stake consensus mechanism and an object-centric data model. Unlike account-based blockchains such as Ethereum, which have a set of rules for transferring ownership of accounts, Sui treats data on the blockchain as programmable objects with defined ownership and transfer rules. This object-oriented approach makes it easier to build asset-tracking apps, making it especially ideal for supply chain use cases \parencite{sui2024documentation}.

Sui's smart contract language is the Move programming language, which has been developed for the Diem blockchain project. The issue of security is also crucial in the design of Move, and its linear type system reduces some frequent safety risks, such as double spending, and reentrancy attacks. The Move-based smart contract introduces an object called \texttt{ProduceBatch} that includes details regarding the crop type, weight, address of the farmer from which the produce was bought, address of the current custodian, and a hash of the data integrity. The smart contract ensures that only the current custodian can transfer the custody, giving a corresponding on-chain access control system that matches the role-based access control in the off-chain backend.

Sui's throughput is provided by parallel transaction execution. If the transactions are related to independent objects, then the network can process them in parallel, without waiting for the consensus of all transactions. This is a very good feature for agricultural applications, which need quick confirmation of batch minting and custody transfers without disrupting the logistics process, hence simple operations benefit from sub-second transaction finality.

\section{Related Work Summary and Comparative Analysis}

Table~\ref{tab:related_work} presents a comparative summary of the key related works reviewed in this chapter, evaluated against the features that TerraNode provides.

\begin{table}[H]
\centering
\small
\singlespacing
\renewcommand{\arraystretch}{1.15}
\caption{Comparative Analysis of Related Works}
\label{tab:related_work}
\begin{tabularx}{\textwidth}{|l|>{\raggedright\arraybackslash}X|c|c|c|c|c|}
\hline
\textbf{Study} & \textbf{Focus Area} & \rotatebox{90}{\textbf{Blockchain}} & \rotatebox{90}{\textbf{Analytics}} & \rotatebox{90}{\textbf{Traceability}} & \rotatebox{90}{\textbf{Web UI}} & \rotatebox{90}{\textbf{Public Ledger}} \\
\hline
\textcite{alabsi2023blockchain} & Smart farm IoT security & Yes & No & No & No & No \\
\hline
\textcite{caro2018blockchain} & Agri-food traceability & Yes & No & Yes & No & Yes \\
\hline
\textcite{khaki2019cnn} & Crop yield prediction & No & Yes & No & No & N/A \\
\hline
\textcite{layek2025secured} & IoT + ML + blockchain & Yes & Yes & No & No & No \\
\hline
\textcite{mirabelli2020blockchain} & Supply chain traceability & Yes & No & Yes & No & Varies \\
\hline
\textcite{ktari2022agricultural} & Embedded blockchain & Yes & No & Yes & No & No \\
\hline
\textcite{sun2025storage} & Storage optimization & Yes & No & Yes & No & Yes \\
\hline
\textcite{cocco2021blockchain} & SSI in food supply chain & Yes & No & Yes & No & Yes \\
\hline
\textcite{garg2018crop} & Fuzzy yield forecasting & No & Yes & No & No & N/A \\
\hline
\textbf{TerraNode} & \textbf{Integrated platform} & \textbf{Yes} & \textbf{Yes} & \textbf{Yes} & \textbf{Yes} & \textbf{Yes} \\
\hline
\end{tabularx}
\end{table}

Comparative analysis shows that some individual elements of the architecture of TerraNode have already been explored in previous studies but there is no system that integrates all the five capabilities: blockchain-based data immutability, predictive analytics, supply chain traceability, role-based web interface and deployment on a public high throughput ledger. Most current blockchain-agriculture systems stress on traceability or security and do not include analytical forecasting. On the other hand, there is little attention paid to data integrity in crop yield prediction studies. The key challenge is addressed by TerraNode, which is a single platform consolidating all of this.

\section{Research Gap}

Through literature review, the following 3 gaps were identified:

Firstly, there is a lack of harmony between data security and data analytics in agriculture. Platforms with blockchain technology are usually more about the act of recording and verifying transactions and offer little to no functionality in terms of analyzing the data being protected by the blockchain. In the meantime, predictive analytics platforms are working on agricultural data, not suspecting that data may have been tampered with. The gap is bridged by making verification of blockchain data an essential step before analytical processing, with only data verifiable through integrity check of its corresponding hash SHA-256 being passed through to the yield prediction engine.

Second, most of the blockchain-agricultural prototypes are built on private and consortium blockchains, which restrict transparency, or public blockchains such as Ethereum, but are too costly to use for farmers' frequent data logging needs. With the Sui blockchain, which is open like a public ledger but has low transaction fees, TerraNode can help overcome this challenge.

Thirdly, current systems usually offer a partial and non-role aware Web interface that does not support all of the needs of the farmers, logistics handlers and administrators in a single system. The frontend provides customisable dashboards for every role, making it easy to understand and accessible to stakeholders of different technical backgrounds.

\section{Conceptual Framework}

This project is conceptualised in the context of two well-known information systems theories. The Technology Acceptance Model (TAM) introduced by Davis states that the acceptance of a new information system is related to perceived usefulness and perceived ease of use. The system is useful for TerraNode, as it generates actionable yield predictions and records transparent supply chain journeys, while being easy to use via an intuitive web-based portal with role-specific dashboards.

The DeLone and McLean Information Systems Success Model identifies six dimensions of system success: information quality, system quality, service quality, user satisfaction, system use, and net benefits. The architecture of TerraNode tackles information quality by ensuring blockchain verification of data integrity, system quality with a modular three-tier structure, and net benefits by delivering measurable results like tamper-proof records and predictive forecasts.

Figure~\ref{fig:conceptual_framework} depicts the conceptual framework, in which the theoretical foundation and the practical system components are linked together.

\begin{figure}[H]
    \centering
    \includegraphics[width=1.0\textwidth]{figures/fig_conceptual_framework.png}
    \caption{Conceptual Framework of the TerraNode System}
    \label{fig:conceptual_framework}
\end{figure}

\section{Chapter Summary}

This chapter examined the various literature studies that investigated the five thematic areas of precision agriculture and digital farming, blockchain technology and its applications in agriculture, smart farm security frameworks, predictive crop yield analytics, and supply chain traceability systems. The review touched upon eighteen published research studies and highlighted their contribution and limitations. The comparative analysis has shown that there is no single system that integrates blockchain-based data immutability, predictive analytics, supply chain custody tracking, a role-based customizable web interface and a public high throughput ledger in one platform. The identified research gaps directly inspire the design and implementation of TerraNode, which is presented in the next chapter.